\makeatletter
\def\input@path{{../styles/}{../../styles/}{../../../styles/}{../}{../../}{../../../}}
\makeatother
\documentclass{ee122_notes}
\input{macros.tex}
\input{preamble.tex}
\renewcommand{\releasedate}{April 27, 2026}
\begin{document}
\section*{EE 122/ME 141 Week 14, Lecture 1: Frequency-domain analysis in presence of noise and disturbances}
\subsection*{Instructor: \instructor}
\subsection*{Date: \releasedate}
\section{Announcements}
\begin{itemize}
    \item Problem set \#11 due on April 29, 2026 
    \item Project final paper due on May 7, 2026 
\end{itemize}
\section{Learning Objectives}
By the end of this lecture, you should be able to:
\begin{itemize}
    \item Understand the effect of noise and disturbances on the performance of a control system in the frequency domain.
    \item Use the frequency-domain analysis tools, such as Bode plots, to analyze the sensitivity of a control system to noise and disturbances.
\end{itemize}

The content of this lecture follows directly from the textbook. So, you are encouraged to read the following sections (in order):
\begin{itemize}
    \item Section 12.1: Sensitivity functions 
    \item Figure 12.1: Explore the block diagram carefully.
    \item Table 12.1: Derivation of all transfer functions.
    \item Example 12.1: Need for more than loop transfer function. (you can stop after this example)
    \item Measuring Specifications on Page 12-12. 
\end{itemize}
\end{document}
